% GENERATED FILE. Edit canonical lessons, then run npm run generate:books.
% canonical-source-hash: fnv1a64:cf100b571e93fca4
% canonical-lessons: PA-C189-want-return, PA-C189-chahunde, PA-C189-chahunde-review, PA-C189-chahunde-han, PA-C189-han-review, PA-C189-chahunde-ho, PA-C189-frame-contrast, PA-C189-payoff

\chapter{They Want; Do You Want?}
\label{ch:pa-a2-want-plural}
\hlchaptermodality{189}

\begin{chapteropening}
\textbf{By the end of this chapter:} \emph{I can recognise and use the shared plural or respectful want form in a familiar statement and question.}

Retrieve the shared plural or respectful want form in a third-person plural statement and a second-person polite question across four skills.
\end{chapteropening}

\section[\texorpdfstring{\pa{ਚਾਹੁਣਾ} · \pa{ਉਹ} \pa{ਕੰਮ} … (chaahunaa; oh …)}{ਚਾਹੁਣਾ · ਉਹ ਕੰਮ … (chaahunaa; oh …)}]{Return to familiar want language}

\label{lesson:PA-C189-want-return}



\begin{quote}
\emph{Hear:} \emph{chaahunaa}; YOU SAY: \textbf{\pa{ਚਾਹੁਣਾ}}
\end{quote}

\begin{grammarlens}[title={Grammar Lens}]
Retrieve the familiar plural or respectful habitual frame before changing the verb: \textbf{\pa{ਉਹ} \pa{ਕੰਮ} \pa{ਕਰਦੇ} \pa{ਹਨ}}.
\end{grammarlens}

\subsection*{Your turn}
\begin{itemize}
\raggedright
  \item \emph{Hear:} \emph{chaahunaa}; YOU SAY: \textbf{\pa{ਚਾਹੁਣਾ}}
  \item \emph{Hear:} \emph{oh kamm karde han}; YOU SAY: \textbf{\pa{ਉਹ} \pa{ਕੰਮ} \pa{ਕਰਦੇ} \pa{ਹਨ}}
\end{itemize}

\begin{tcolorbox}[breakable,colback=teal!4,colframe=teal!35!black,title={Before you move on}]
Say the infinitive, then the familiar agreement frame. This opener owns no new form.
\end{tcolorbox}
\section[\texorpdfstring{\pa{ਚਾਹੁੰਦੇ} (chaahunde)}{ਚਾਹੁੰਦੇ (chaahunde)}]{The plural or respectful want form}

\label{lesson:PA-C189-chahunde}



\begin{quote}
Retrieve \textbf{\pa{ਚਾਹੁਣਾ}}, then the middle form in \textbf{\pa{ਉਹ} \pa{ਕੰਮ} \pa{ਕਰਦੇ} \pa{ਹਨ}}.
\end{quote}

\begin{grammarlens}[title={Grammar Lens: direct owner}]
\begin{quote}
\textbf{\pa{ਚਾਹੁੰਦੇ}} — \emph{chaahunde} — want, with plural or respectful habitual agreement
\end{quote}

Treat the whole written form as today’s measured target. The next lesson retrieves it before either complete sentence frame is built.
\end{grammarlens}

\subsection*{Your turn}
\begin{itemize}
\raggedright
  \item \emph{Hear:} \emph{chaahunde}; YOU SAY: \textbf{\pa{ਚਾਹੁੰਦੇ}}
  \item \emph{Read it:} \textbf{\pa{ਚਾਹੁੰਦੇ}}
  \item \emph{Write it:} copy \textbf{\pa{ਚਾਹੁੰਦੇ}} once
\end{itemize}

\begin{tcolorbox}[breakable,colback=teal!4,colframe=teal!35!black,title={Before you move on}]
Produce the plural or respectful want form without looking.
\end{tcolorbox}
\section[\texorpdfstring{\pa{ਚਾਹੁੰਦੇ} (chaahunde)}{ਚਾਹੁੰਦੇ (chaahunde)}]{Retrieve the new agreement form}

\label{lesson:PA-C189-chahunde-review}



\begin{quote}
Hear “want, plural or respectful” and prepare the single form.
\end{quote}

\begin{grammarlens}[title={Grammar Lens}]
Keep \textbf{\pa{ਚਾਹੁੰਦੇ}} distinct from the infinitive \textbf{\pa{ਚਾਹੁਣਾ}}. Its agreement job matches the familiar plural or respectful pattern in \textbf{\pa{ਕਰਦੇ}}.
\end{grammarlens}

\subsection*{Your turn}
\begin{itemize}
\raggedright
  \item \emph{Hear:} \emph{chaahunde}; YOU SAY: \textbf{\pa{ਚਾਹੁੰਦੇ}}
  \item \emph{Read it:} \textbf{\pa{ਚਾਹੁੰਦੇ}}
  \item \emph{Write it:} cover the model and write \textbf{\pa{ਚਾਹੁੰਦੇ}}
\end{itemize}

\begin{tcolorbox}[breakable,colback=teal!4,colframe=teal!35!black,title={Before you move on}]
Read the form once more and state its bounded agreement job.
\end{tcolorbox}
\section[\texorpdfstring{\pa{ਉਹ} \pa{ਕੰਮ} \pa{ਕਰਨਾ} \pa{ਚਾਹੁੰਦੇ} \pa{ਹਨ} (oh kamm karnaa …)}{ਉਹ ਕੰਮ ਕਰਨਾ ਚਾਹੁੰਦੇ ਹਨ (oh kamm karnaa …)}]{Build the third-person plural frame}

\label{lesson:PA-C189-chahunde-han}



\begin{quote}
Retrieve \textbf{\pa{ਉਹ}}, \textbf{\pa{ਕੰਮ}}, \textbf{\pa{ਕਰਨਾ}}, and \textbf{\pa{ਚਾਹੁੰਦੇ}}.
\end{quote}

\begin{grammarlens}[title={Grammar Lens: complete statement}]
\begin{quote}
\textbf{\pa{ਉਹ} \pa{ਕੰਮ} \pa{ਕਰਨਾ} \pa{ਚਾਹੁੰਦੇ} \pa{ਹਨ}} — \emph{oh kamm karnaa chaahunde han} — they want to work
\end{quote}

Read the plural subject, the wanted action, the agreement form, and the familiar plural auxiliary as one bounded frame.
\end{grammarlens}

\subsection*{Your turn}
\begin{itemize}
\raggedright
  \item \emph{Hear:} \emph{oh kamm karnaa chaahunde han}; YOU SAY: \textbf{\pa{ਉਹ} \pa{ਕੰਮ} \pa{ਕਰਨਾ} \pa{ਚਾਹੁੰਦੇ} \pa{ਹਨ}}
  \item \emph{Read it:} \textbf{\pa{ਉਹ} \pa{ਕੰਮ} \pa{ਕਰਨਾ} \pa{ਚਾਹੁੰਦੇ} \pa{ਹਨ}}
  \item \emph{Write it:} copy \textbf{\pa{ਚਾਹੁੰਦੇ} \pa{ਹਨ}} once
\end{itemize}

\begin{tcolorbox}[breakable,colback=teal!4,colframe=teal!35!black,title={Before you move on}]
Produce the complete “they want to work” frame.
\end{tcolorbox}
\section[\texorpdfstring{\pa{ਉਹ} \pa{ਕੰਮ} \pa{ਕਰਨਾ} \pa{ਚਾਹੁੰਦੇ} \pa{ਹਨ} (oh kamm karnaa …)}{ਉਹ ਕੰਮ ਕਰਨਾ ਚਾਹੁੰਦੇ ਹਨ (oh kamm karnaa …)}]{Retrieve the plural statement}

\label{lesson:PA-C189-han-review}



\begin{quote}
Hear “they want to work” and prepare the whole statement.
\end{quote}

\begin{grammarlens}[title={Grammar Lens}]
Keep the target next to its familiar plural auxiliary. The frame remains \textbf{\pa{ਉਹ} \pa{ਕੰਮ} \pa{ਕਰਨਾ} \pa{ਚਾਹੁੰਦੇ} \pa{ਹਨ}}.
\end{grammarlens}

\subsection*{Your turn}
\begin{itemize}
\raggedright
  \item \emph{Hear:} \emph{oh kamm karnaa chaahunde han}; YOU SAY: \textbf{\pa{ਉਹ} \pa{ਕੰਮ} \pa{ਕਰਨਾ} \pa{ਚਾਹੁੰਦੇ} \pa{ਹਨ}}
  \item \emph{Read it:} \textbf{\pa{ਉਹ} \pa{ਕੰਮ} \pa{ਕਰਨਾ} \pa{ਚਾਹੁੰਦੇ} \pa{ਹਨ}}
  \item \emph{Write it:} cover and write \textbf{\pa{ਉਹ} \pa{ਕੰਮ} \pa{ਕਰਨਾ} \pa{ਚਾਹੁੰਦੇ} \pa{ਹਨ}}
\end{itemize}

\begin{tcolorbox}[breakable,colback=teal!4,colframe=teal!35!black,title={Before you move on}]
Say the statement once more without looking.
\end{tcolorbox}
\section[\texorpdfstring{\pa{ਤੁਸੀਂ} \pa{ਕੀ} \pa{ਕਰਨਾ} \pa{ਚਾਹੁੰਦੇ} \pa{ਹੋ} (tusiin kii …)}{ਤੁਸੀਂ ਕੀ ਕਰਨਾ ਚਾਹੁੰਦੇ ਹੋ (tusiin kii …)}]{Ask what someone wants to do}

\label{lesson:PA-C189-chahunde-ho}



\begin{quote}
Retrieve \textbf{\pa{ਤੁਸੀਂ}}, \textbf{\pa{ਕੀ}}, \textbf{\pa{ਕਰਨਾ}}, \textbf{\pa{ਚਾਹੁੰਦੇ}}, and \textbf{\pa{ਹੋ}}.
\end{quote}

\begin{grammarlens}[title={Grammar Lens: complete question}]
\begin{quote}
\textbf{\pa{ਤੁਸੀਂ} \pa{ਕੀ} \pa{ਕਰਨਾ} \pa{ਚਾਹੁੰਦੇ} \pa{ਹੋ}} — \emph{tusiin kii karnaa chaahunde ho} — what do you want to do?
\end{quote}

The respectful or plural subject keeps the same target form, while the familiar second-person auxiliary closes the question.
\end{grammarlens}

\subsection*{Your turn}
\begin{itemize}
\raggedright
  \item \emph{Hear:} \emph{tusiin kii karnaa chaahunde ho}; YOU SAY: \textbf{\pa{ਤੁਸੀਂ} \pa{ਕੀ} \pa{ਕਰਨਾ} \pa{ਚਾਹੁੰਦੇ} \pa{ਹੋ}}
  \item \emph{Read it:} \textbf{\pa{ਤੁਸੀਂ} \pa{ਕੀ} \pa{ਕਰਨਾ} \pa{ਚਾਹੁੰਦੇ} \pa{ਹੋ}}
  \item \emph{Write it:} copy \textbf{\pa{ਚਾਹੁੰਦੇ} \pa{ਹੋ}} once
\end{itemize}

\begin{tcolorbox}[breakable,colback=teal!4,colframe=teal!35!black,title={Before you move on}]
Ask the complete question without looking.
\end{tcolorbox}
\section[\texorpdfstring{\pa{ਉਹ} \pa{ਕੰਮ} \pa{ਕਰਨਾ} \pa{ਚਾਹੁੰਦੇ} … (oh kamm karnaa …)}{ਉਹ ਕੰਮ ਕਰਨਾ ਚਾਹੁੰਦੇ … (oh kamm karnaa …)}]{Contrast the two complete frames}

\label{lesson:PA-C189-frame-contrast}



\begin{quote}
Prepare the plural statement and the respectful question.
\end{quote}

\begin{grammarlens}[title={Grammar Lens: same target, different frame endings}]
\begin{itemize}
\raggedright
  \item \textbf{\pa{ਉਹ} \pa{ਕੰਮ} \pa{ਕਰਨਾ} \pa{ਚਾਹੁੰਦੇ} \pa{ਹਨ}} — a third-person plural statement.
  \item \textbf{\pa{ਤੁਸੀਂ} \pa{ਕੀ} \pa{ਕਰਨਾ} \pa{ਚਾਹੁੰਦੇ} \pa{ਹੋ}} — a second-person plural or respectful question.
\end{itemize}
\end{grammarlens}

\subsection*{Your turn}
\begin{itemize}
\raggedright
  \item \emph{Say it:} \textbf{\pa{ਉਹ} \pa{ਕੰਮ} \pa{ਕਰਨਾ} \pa{ਚਾਹੁੰਦੇ} \pa{ਹਨ}}
  \item \emph{Say it:} \textbf{\pa{ਤੁਸੀਂ} \pa{ਕੀ} \pa{ਕਰਨਾ} \pa{ਚਾਹੁੰਦੇ} \pa{ਹੋ}}
  \item \emph{Write it:} \textbf{\pa{ਚਾਹੁੰਦੇ} \pa{ਹਨ} · \pa{ਚਾਹੁੰਦੇ} \pa{ਹੋ}}
\end{itemize}

\begin{tcolorbox}[breakable,colback=teal!4,colframe=teal!35!black,title={Before you move on}]
Cover the models and produce statement, then question.
\end{tcolorbox}
\section[\texorpdfstring{\pa{ਉਹ} \pa{ਕੰਮ} \pa{ਕਰਨਾ} \pa{ਚਾਹੁੰਦੇ} … (oh kamm karnaa …)}{ਉਹ ਕੰਮ ਕਰਨਾ ਚਾਹੁੰਦੇ … (oh kamm karnaa …)}]{A bounded plural and respectful want payoff}

\label{lesson:PA-C189-payoff}



\begin{quote}
Recall the target alone, then prepare both complete frames.
\end{quote}

\subsection*{Your turn: four skills}
\begin{itemize}
\raggedright
  \item \emph{Hear:} \emph{oh kamm karnaa chaahunde han; tusiin kii karnaa chaahunde ho}; YOU SAY: \textbf{\pa{ਉਹ} \pa{ਕੰਮ} \pa{ਕਰਨਾ} \pa{ਚਾਹੁੰਦੇ} \pa{ਹਨ}; \pa{ਤੁਸੀਂ} \pa{ਕੀ} \pa{ਕਰਨਾ} \pa{ਚਾਹੁੰਦੇ} \pa{ਹੋ}}
  \item \emph{Read it:} \textbf{\pa{ਉਹ} \pa{ਕੰਮ} \pa{ਕਰਨਾ} \pa{ਚਾਹੁੰਦੇ} \pa{ਹਨ} · \pa{ਤੁਸੀਂ} \pa{ਕੀ} \pa{ਕਰਨਾ} \pa{ਚਾਹੁੰਦੇ} \pa{ਹੋ}}
  \item \emph{Write it:} cover the line and write \textbf{\pa{ਚਾਹੁੰਦੇ} \pa{ਹਨ} · \pa{ਚਾਹੁੰਦੇ} \pa{ਹੋ}}
\end{itemize}

\begin{tcolorbox}[breakable,colback=teal!4,colframe=teal!35!black,title={Before you move on}]
Produce the plural statement and the respectful question once more.

This is the twelfth prerequisite tranche exposed by Punjabi A2 mock drafting. It does not unblock Punjabi A2 mock 1 by itself. It is not an A2 mock, not A2 readiness or attainment, and not external or human validation.
\end{tcolorbox}
